\documentclass[a4paper,11pt]{article}

\usepackage{fullpage}
\usepackage[utf8]{inputenc}
\usepackage[T1]{fontenc}
\usepackage{amsmath,amsthm}
\usepackage{amsfonts}
\usepackage{graphicx}
\usepackage{latexsym}
\usepackage{float}
\usepackage{comment}
\usepackage{array}
\usepackage{booktabs}

% --- Packages added for pedagogical enrichment ---
\usepackage{xcolor}
\usepackage{listings}
\usepackage[most]{tcolorbox}
\usepackage{hyperref}

%%%%%%%%%%%%%%% LISTINGS CONFIGURATION %%%%%%%%%%%%%%%%%%%

\definecolor{codegreen}{rgb}{0.25,0.49,0.25}
\definecolor{codegray}{rgb}{0.5,0.5,0.5}
\definecolor{codepurple}{rgb}{0.58,0,0.82}
\definecolor{codered}{rgb}{0.7,0.1,0.1}
\definecolor{codeblue}{rgb}{0.0,0.0,0.7}
\definecolor{backcolour}{rgb}{0.97,0.97,0.97}

\lstset{
    language=C,
    backgroundcolor=\color{backcolour},
    commentstyle=\color{codegreen}\itshape,
    keywordstyle=\color{codeblue}\bfseries,
    stringstyle=\color{codered},
    basicstyle=\ttfamily\small,
    breakatwhitespace=false,
    breaklines=true,
    captionpos=b,
    keepspaces=true,
    numbers=left,
    numbersep=8pt,
    numberstyle=\tiny\color{codegray},
    showspaces=false,
    showstringspaces=false,
    showtabs=false,
    tabsize=4,
    frame=single,
    rulecolor=\color{codegray},
    xleftmargin=1.5em,
    framexleftmargin=1.5em,
    literate={é}{{\'e}}1 {è}{{\`e}}1 {ê}{{\^e}}1 {à}{{\`a}}1 {ù}{{\`u}}1 {î}{{\^i}}1 {ô}{{\^o}}1
}

%%%%%%%%%%%%%%% TCOLORBOX ENVIRONMENTS %%%%%%%%%%%%%%%%%%

% Course reminder (light blue background)
\newtcolorbox{rappel}[1][]{
    colback=blue!5!white,
    colframe=blue!50!black,
    fonttitle=\bfseries,
    title={\large Course Reminder},
    #1
}

% Key takeaway (light green background)
\newtcolorbox{pointcle}[1][]{
    colback=green!5!white,
    colframe=green!40!black,
    fonttitle=\bfseries,
    title={Key Takeaway},
    #1
}

% Warning (light orange background)
\newtcolorbox{attention}[1][]{
    colback=orange!8!white,
    colframe=orange!60!black,
    fonttitle=\bfseries,
    title={Warning},
    #1
}

% Ethical disclaimer (red background)
\newtcolorbox{ethique}[1][]{
    colback=red!5!white,
    colframe=red!60!black,
    fonttitle=\bfseries\large,
    title={Ethical and Legal Disclaimer},
    #1
}

%%%%%%%%%%%%%%% COMMANDS %%%%%%%%%%%%%%%%%%%%%%%%%%%

\newcommand{\itemb}{\item[$\bullet$]}

\newcommand{\Fd}{\mathbb{F}}
\newcommand{\Znat}{\mathbb{Z}}
\newcommand{\Nnat}{\mathbb{N}}

%%%%%%%%%%%%%%% ENVIRONMENTS %%%%%%%%%%%%%%%%%%%%%%%

\theoremstyle{plain}
\newtheorem{lemme}{Lemma}
\newtheorem{algorithme}{Algorithm}
\newtheorem{corolaire}{Corollary}
\newtheorem{propriete}{Property}
\newtheorem{proposition}{Proposition}
\newtheorem{theoreme}{Theorem}
\newtheorem{definition}{Definition}

\theoremstyle{definition}
\newtheorem{exercice}{Exercise}
\newtheorem{remarque}{Remark}
\newtheorem{exemple}{Example}

\hypersetup{
    colorlinks=true,
    linkcolor=blue!60!black,
    urlcolor=blue!60!black
}

%%%%%%%%%%%%%% DOCUMENT %%%%%%%%%%%%%%%%%%%%%%%%%%%%


\begin{document}

\hbox{\raisebox{0.4em}{\vrule depth 0pt height 0.5pt width \textwidth}}
\noindent \textbf{University of Perpignan} \hfill  L3 Computer Science \\
 \hfill C. Legrand \hfill  Year \the\year \\
\smallskip
\begin{center}
{
  {\scshape \large Lab Work 5 -- Security}  \\
 Intrusion Techniques -- Buffer Overflow
}
\end{center}
\hbox{\raisebox{0.4em}{\vrule depth 0pt height 0.9pt width \textwidth}}

\bigskip

% --- Ethical disclaimer ---
\begin{ethique}
The techniques presented in this lab work are studied within a \textbf{strictly academic context} in order to understand the mechanisms used by attackers and thus be better able to defend against them.

\textbf{Any use of these techniques outside the educational context is illegal} and punishable by criminal penalties (Chapter III of the French Penal Code, Articles 323-1 to 323-8 -- \textit{Offences against Automated Data Processing Systems}):
\begin{itemize}
    \item Unauthorized access to or remaining in an ADPS (Art.~323-1): \textbf{3 years} imprisonment and \textbf{\texteuro{}100,000} fine
    \item Obstructing the operation of an ADPS (Art.~323-2): \textbf{5 years} imprisonment and \textbf{\texteuro{}150,000} fine
    \item Fraudulent introduction of data into an ADPS (Art.~323-3): \textbf{5 years} imprisonment and \textbf{\texteuro{}150,000} fine
\end{itemize}
Penalties increased to \textbf{7 years} and \textbf{\texteuro{}300,000} when the targeted system processes personal data operated by the State, and up to \textbf{10 years} and \textbf{\texteuro{}300,000} for organized crime (Art.~323-4-1). \textit{Penalties updated by the LOPMI law no.~2023-22 of January 24, 2023.}
\end{ethique}

\bigskip

\begin{tcolorbox}[colback=gray!5!white, colframe=gray!60!black, title={Compilation Options}]
The following options should be used when you want to overflow the stack (\texttt{-fno-stack-protector}) and know the code addresses (\texttt{-no-pie}) which are no longer randomized:
\begin{lstlisting}[language=bash,numbers=none]
gcc -fno-stack-protector -no-pie -g code.c
\end{lstlisting}
One last option that we will not use: if you wish to execute code injected into the stack (\texttt{-z execstack}).
\end{tcolorbox}

\medskip

\begin{rappel}
\textbf{Stack organization:} when a function is called, the processor pushes:
\begin{enumerate}
    \item The function parameters (partially via registers in x86-64)
    \item The return address \texttt{savedRIP} (pushed by \texttt{call})
    \item The old \texttt{RBP} (\texttt{savedRBP}) saved by \texttt{push rbp}
    \item The local variables of the function
\end{enumerate}
The \texttt{RSP} register points to the top of the stack; \texttt{RBP} points to the base of the current stack frame.
\end{rappel}

\bigskip


%%%%%%%%%%%%%%%%%%%%%%%%%%%%%%%%%%%%%%%%%%%%%%%%%%%%%%%%%%%%%%%%%%
%% EXERCISE 1: Simple Buffer Overflow
%%%%%%%%%%%%%%%%%%%%%%%%%%%%%%%%%%%%%%%%%%%%%%%%%%%%%%%%%%%%%%%%%%

\begin{exercice}[Buffer Overflow]
\begin{enumerate}
\item
Consider the following program:
\begin{lstlisting}[language=C]
#include <stdio.h>
#include <string.h>

int main(int argc, char *argv[]){
  int valeur=5;
  char tampon_un[8], tampon_deux[8];

  strcpy(tampon_un,"un");
  strcpy(tampon_deux,"deux");

  printf("[AVANT] tampon_deux est a %p et contient '%s'\n",
         tampon_deux, tampon_deux);
  printf("[AVANT] tampon_un est a %p et contient '%s'\n",
         tampon_un, tampon_un);
  printf("[AVANT] valeur est a %p et vaut %d (0x%08x)\n",
         &valeur, valeur, valeur);

  printf("\n [STRCPY] copie de %d octets dans tampon_deux\n\n",
         (int)strlen(argv[1]));
  strcpy(tampon_deux, argv[1]);

  printf("[APRES] tampon_deux est a %p et contient '%s'\n",
         tampon_deux, tampon_deux);
  printf("[APRES] tampon_un est a %p et contient '%s'\n",
         tampon_un, tampon_un);
  printf("[APRES] valeur est a %p et vaut %d (0x%08x)\n",
         &valeur, valeur, valeur);

  return(0);
}
\end{lstlisting}

By varying the argument, trigger a buffer overflow and analyze it with \texttt{gdb}.

\begin{pointcle}
The \texttt{strcpy} function does \textbf{not} check the size of the destination buffer. If the source is longer than the buffer, the extra bytes overwrite adjacent data on the stack.
\end{pointcle}

\item  Now consider the following code snippet which restricts access rights:

\begin{lstlisting}[language=C]
#include <stdio.h>
#include <stdlib.h>
#include <string.h>

int verif_authentication(char *password){
  int auth_flag = 0;
  char password_tampon[16];

  strcpy(password_tampon, password);

  if( strcmp(password_tampon, "brillig") == 0)
    auth_flag = 1;
  if(strcmp(password_tampon, "outgrabe") == 0)
    auth_flag = 1;

  return auth_flag;
}

int main( int argc, char *argv[]){
  if(argc < 2){
    printf("Usage: %s <password>\n", argv[0]);
    exit(0);
  }

  if(verif_authentication(argv[1])){
    printf("\n-=-=-=-=-=-=-=-=-=-=--=-=-\n");
    printf("      Acces accorde.\n");
    printf("-=-=-=-=-=-=-=-=-=-=--=-=-\n");
  }
  else
    printf("     Acces refuse\n");
}
\end{lstlisting}
Is it possible to perform a buffer overflow to force access (\texttt{auth\_flag} $\neq 0$)? What happens if the variables are declared in the following order:
\begin{lstlisting}[language=C,numbers=none]
  char password_buffer[16];
  int auth_flag = 0;
\end{lstlisting}
Can you work around this issue to still display \texttt{Acces accorde} via buffer overflow? Explain in detail why it works.
\end{enumerate}
\end{exercice}


%%%%%%%%%%%%%%%%%%%%%%%%%%%%%%%%%%%%%%%%%%%%%%%%%%%%%%%%%%%%%%%%%%
%% EXERCISE 2: Return Address Modification (Increment)
%%%%%%%%%%%%%%%%%%%%%%%%%%%%%%%%%%%%%%%%%%%%%%%%%%%%%%%%%%%%%%%%%%

\begin{exercice}[Return Address Modification on the Stack]
Take the code seen in the lecture (given below) and adjust the values of $a$ and $b$ in order to execute upon function return:
\begin{itemize}
\item Directly the end of the \texttt{main} function (thus skipping both \texttt{printf("coucou1")} and \texttt{printf("coucou2")}).
  \item What will happen if we return to the instruction \texttt{b=15}?
\end{itemize}

\begin{lstlisting}[language=C]
#include <stdio.h>

void function(long a, long b)
{
  long buffer[2]={0,1};
  buffer[a]+=b;
}
void main()
{
  long a,b;
  a=1;
  b=15;
  function(a,b);
  printf("coucou1\n");
  printf("coucou2\n");
}
\end{lstlisting}

\begin{rappel}
To find the correct values of \texttt{a} and \texttt{b}, you need to:
\begin{enumerate}
    \item Analyze the assembly code with \texttt{gdb} to determine the size of the stack frame of \texttt{function}
    \item Calculate the offset between \texttt{buffer[0]} and \texttt{savedRIP}
    \item Determine the necessary increment to jump over the desired instructions
\end{enumerate}
\end{rappel}
\end{exercice}


%%%%%%%%%%%%%%%%%%%%%%%%%%%%%%%%%%%%%%%%%%%%%%%%%%%%%%%%%%%%%%%%%%
%% EXERCISE 3: Return Address Overwrite (Overflow)
%%%%%%%%%%%%%%%%%%%%%%%%%%%%%%%%%%%%%%%%%%%%%%%%%%%%%%%%%%%%%%%%%%

\begin{exercice}[Return Address Overwrite via Overflow]
Consider the following code:
\begin{lstlisting}[language=C]
#include <stdio.h>
#include <stdlib.h>
#include <string.h>

void f(char *s)
{
  char nom[10];
  strcpy(nom,s);
  printf("\n");
  printf("%lx\n",*((unsigned long *)(nom+18)));
  printf("Bonjour %s, on est dans f\n",nom);
}

void secret()
{
  printf("on est dans secret\n");
}

int main(int argc, char **argv)
{
  printf("adresse de f=%p, et de secret=%p\n",f,secret);
  f(argv[1]);
}
\end{lstlisting}
By overwriting the return address of the call to function \texttt{f}, make it so that the function \texttt{secret} is executed.

\begin{attention}
To pass a string of characters specified by their ASCII code as a parameter, you can use \texttt{python3} as follows (for the two ASCII codes \texttt{0x11} and \texttt{0xff}):
\begin{lstlisting}[language=bash,numbers=none]
$ python3 -c "print('\x11\xff')"
\end{lstlisting}
\end{attention}
\end{exercice}


%%%%%%%%%%%%%%%%%%%%%%%%%%%%%%%%%%%%%%%%%%%%%%%%%%%%%%%%%%%%%%%%%%
%% EXERCISE 4: printf Format String
%%%%%%%%%%%%%%%%%%%%%%%%%%%%%%%%%%%%%%%%%%%%%%%%%%%%%%%%%%%%%%%%%%

\begin{exercice}[Exploiting the \texttt{printf} Format String]
Consider the following code compiled with the options \texttt{-fno-stack-protector -no-pie}:
\begin{lstlisting}[language=C]
#include <stdio.h>
#include <stdlib.h>
#include <string.h>

int main(int argc, char *argv[]) {
  unsigned long d;
  char texte[512];

  static int test_var = 100*256*256+100*256+100;
  if(argc < 3) {
    printf("Usage: %s <text to print> <entier>\n", argv[0]);
    exit(0);
  }
  d=atoi(argv[2]);

  strcpy(texte, argv[1]);
  printf("La bonne facon d'afficher des entrees utilisateurs:\n");
  printf("%s", texte);
  printf("\nLa mauvaise facon d'afficher des entrees utilisateurs:\n");
  printf(texte);
  printf("\n");
  // Debug output
  printf("[*] test_var @ %p = %lu, %04x\n", &test_var, &test_var, test_var);
  printf("[*] adresse de d=%p, valeur de d=0x%016lx\n", &d, d);
  printf("[*] adresse de texte=%p\n", texte);
  exit(0);
}
\end{lstlisting}

\begin{enumerate}
\item Understand and test the program, then apply what was seen in the lecture to display the contents of the \texttt{main} stack by carefully choosing the program argument on the command line.
\item If the program argument contains the format \texttt{"\%s"}, what will happen? Can you choose the address that \texttt{"\%s"} refers to and use it to read the value of \texttt{test\_var}?
\end{enumerate}

\begin{pointcle}
The call \texttt{printf(texte)} is dangerous because the attacker controls the format string. They can use \texttt{\%x}, \texttt{\%lx}, \texttt{\%s} to read memory, and even \texttt{\%n} to write to memory!
\end{pointcle}
\end{exercice}


\end{document}
